% 中文演示文稿模板：ctexbeamer（beamer + ctex），16:9，XeLaTeX，字体集 fandol（离线可用）。
% 不含参考文献：引用来源写在对应页的"来源："一行或脚注中。
% 结构：标题页 → 目录 → 研究背景 → 方法（公式）→ 结果（三线表）→ 总结 → 致谢页。
% 建议每页只讲一个要点；页数不等于演讲时长。
\documentclass[UTF8,fontset=fandol,aspectratio=169]{ctexbeamer}

\usetheme{Boadilla}
\setbeamertemplate{navigation symbols}{}
\usepackage{amsmath,amssymb}
\usepackage{booktabs}
\usepackage{graphicx}

\title{报告题目}
\subtitle{副标题（可选）}
\author{报告人}
\institute{单位名称}
% 写明会议与日期；\today 会显示该版本内容首次构建的日期（构建时钟按快照固定，保证可复现）。
\date{会议或讲座名称，年月}

\begin{document}

\begin{frame}
  \titlepage
\end{frame}

\begin{frame}{目录}
  \tableofcontents
\end{frame}

\section{研究背景}
\begin{frame}{研究背景}
  % 要解决的问题、为什么值得关注、听众需要记住的一个核心观点。
  \begin{itemize}
    \item 本报告要解决的问题
    \item 这一问题的重要性
    \item 希望听众记住的核心观点
  \end{itemize}
\end{frame}

\section{方法}
\begin{frame}{核心思路}
  % 替换为核心公式，并在本页解释每个符号。
  \begin{equation}
    \hat{\theta} = \operatorname*{arg\,min}_{\theta}
      \sum_{i=1}^{n} \ell\bigl(f_{\theta}(x_i),\, y_i\bigr)
  \end{equation}
  \begin{itemize}
    \item $f_{\theta}$：参数为 $\theta$ 的模型
    \item $\ell$：衡量预测与目标差异的损失函数
  \end{itemize}
\end{frame}

\section{实验结果}
\begin{frame}{实验结果}
  % 只报告真实测得的数据，并注明来源。
  \begin{table}
    \centering
    \begin{tabular}{lcc}
      \toprule
      方法     & 指标一 & 指标二 \\
      \midrule
      基线方法 & --     & --     \\
      本文方法 & --     & --     \\
      \bottomrule
    \end{tabular}
    \caption{表头为示例，数据待填}
  \end{table}
\end{frame}

\section{总结}
\begin{frame}{总结}
  \begin{enumerate}
    \item 第一条结论
    \item 第二条结论
    \item 局限与下一步工作
  \end{enumerate}
\end{frame}

\begin{frame}
  \centering
  \Large 谢谢！欢迎提问。
\end{frame}

\end{document}
